\documentclass[11pt,a4paper]{article}
\usepackage[T1]{fontenc}
\usepackage[utf8]{inputenc}
\usepackage[main=italian,provide=*]{babel}
\usepackage{amsmath,amssymb}
\usepackage{geometry}
\geometry{margin=2.3cm,top=2.5cm,bottom=2.7cm}
\usepackage{booktabs}
\usepackage[table,dvipsnames]{xcolor}
\usepackage{tcolorbox}
\tcbuselibrary{skins,breakable}
\usepackage{titlesec}
\usepackage{enumitem}
\usepackage{needspace}
\usepackage{fancyhdr}
\usepackage{graphicx}
\usepackage{hyperref}
\hypersetup{colorlinks=true,linkcolor=Sepia,citecolor=Sepia,urlcolor=Sepia}
\newcommand{\Tr}{\operatorname{Tr}}
\newcommand{\ket}[1]{\lvert #1 \rangle}
\newcommand{\bra}[1]{\langle #1 \rvert}

% --- palette ---
\definecolor{inkblue}{HTML}{1B3A5C}
\definecolor{lightblue}{HTML}{EAF1F8}
\definecolor{accent}{HTML}{B0562A}
\definecolor{softgray}{HTML}{6B6B6B}
\definecolor{okgreen}{HTML}{2E6B3E}
\definecolor{okgreenbg}{HTML}{EAF5EC}
\definecolor{warnred}{HTML}{9E3B3B}
\definecolor{warnredbg}{HTML}{FBEDED}

% --- titoli di sezione ---
\titleformat{\section}
  {\normalfont\Large\bfseries\color{inkblue}}
  {\thesection}{0.6em}{}[{\color{inkblue!40}\titlerule}]
\titlespacing{\section}{0pt}{0.9em}{0.4em}
\titleformat{\subsection}
  {\normalfont\large\bfseries\color{accent}}
  {}{0em}{}
\titlespacing{\subsection}{0pt}{0.6em}{0.2em}

\pagestyle{fancy}
\fancyhf{}
\renewcommand{\headrulewidth}{0.4pt}
\fancyhead[L]{\small\color{softgray}Test di readout asimmetrico --- correlatori dinamici sotto rumore}
\fancyhead[R]{\small\color{softgray}\thepage}
\fancyfoot[C]{\small\color{softgray}Tesi triennale, Samuele --- Relatore: Prof.\ Alessandro Chiesa}

% --- box riutilizzabili ---
\newtcolorbox{keyeq}[1][]{
  colback=lightblue, colframe=inkblue!70, boxrule=0.6pt,
  arc=2pt, left=8pt, right=8pt, top=3pt, bottom=3pt, #1
}
\newtcolorbox{okbox}[1]{
  colback=okgreenbg, colframe=okgreen!60, boxrule=0.6pt,
  arc=2pt, left=7pt, right=7pt, top=2pt, bottom=2pt,
  fonttitle=\bfseries\color{okgreen}, title={#1}
}
\newtcolorbox{warnbox}[1]{
  colback=warnredbg, colframe=warnred!60, boxrule=0.6pt,
  arc=2pt, left=7pt, right=7pt, top=2pt, bottom=2pt,
  fonttitle=\bfseries\color{warnred}, title={#1}
}

\title{\vspace{-1.5em}\color{inkblue}Test di readout asimmetrico\\[0.2em]
\Large\normalfont\color{softgray}Correlatori dinamici sotto rumore --- processo, risultati, considerazioni}
\author{Note di lavoro --- Tesi triennale, Samuele\\ Relatore: Prof.\ Alessandro Chiesa}
\date{}

\begin{document}
\sloppy
\maketitle
\thispagestyle{fancy}
\vspace{-2.0em}

File di codice: \texttt{correlatori\_readout\_asimmetrico.py},
\texttt{validate\_\allowbreak correlatori\_readout\_asimmetrico.py}. Punto di lavoro:
$b/J=0.35$, $D/J=0.80$, $J=1$, $t=2$ (stesso correlatore di riferimento
del Passo 4, $C_{21}^{xx}(t)$). Teoria completa in
\texttt{teoria\_readout\_asimmetrico.tex}.

\tableofcontents

\section{Obiettivo}

Il relatore ha chiesto esplicitamente (\texttt{domande\_relatore.md},
sez.~12): \emph{``Sull'errore di misura, parti simmetrico. Poi vedi se hai
tempo di fare un test asimmetrico, ma senza perderci troppo tempo''}. Il
readout simmetrico è già stato usato ovunque nella pipeline, con il
risultato che $N^*=5$ (correlatori dinamici, Passo 4) è dimostrato
indipendente da $p_\text{readout}$ per qualunque valore simmetrico. Questo
documento affronta la parte lasciata esplicitamente aperta: cosa succede
quando il readout non è più simmetrico.

La teoria (\texttt{teoria\_readout\_asimmetrico.tex}) ha già mostrato che
l'argomento di invarianza usato per il caso simmetrico --- una pura
riscalatura moltiplicativa del correlatore --- \textbf{non si estende
automaticamente} al caso asimmetrico, perché introduce anche un offset
additivo costante nel piano complesso che non commuta con l'operazione di
modulo. La domanda del relatore ha quindi una risposta aperta a priori,
non scontata: questo documento la chiude numericamente.

\section{Metodologia}

\subsection{Formula applicata}

Per ogni $N$, il correlatore complesso rumoroso viene corretto con la
formula generale derivata nella teoria,
\begin{keyeq}
\centering
$C_\text{readout}(N) = \alpha\,C_\text{ideale}(N) + \beta(1+i), \qquad
\alpha = 1-p_{01}-p_{10}, \qquad \beta = p_{10}-p_{01}$
\end{keyeq}
al posto della correzione simmetrica $(1-2p)$ già in uso. Il limite
$p_{01}=p_{10}=p$ deve riprodurre esattamente il ramo simmetrico --- primo
controllo, non un'assunzione.

\subsection{Ricerca di valori reali: cosa è stato trovato}
\label{sec:ricerca}

La fonte di calibrazione già citata in tesi per \texttt{ibm\_torino}
(Stenger \emph{et al.}, arXiv:2504.15187) riporta solo la mediana
simmetrica $p_\text{readout}=2.3\times10^{-2}$, non la scomposizione in
$p_{01}$ (\texttt{prob\_meas1\_prep0}) e $p_{10}$
(\texttt{prob\_meas0\_prep1}). È stata condotta una ricerca esplicita di
uno split realmente registrato per \texttt{ibm\_torino}: un solo paper
(AbuGhanem, arXiv:2410.00916) presenta una tabella dedicata a
\texttt{ibm\_torino} con i valori per-qubit, ma pubblicata come immagine
non estraibile in forma testuale citabile con i numeri esatti. Nessun'
altra fonte citabile riporta lo split per questo chip specifico.

La ricerca ha invece trovato split reali, citabili, su \textbf{altri} chip
IBM comparabili (stessa tecnologia, stesso ordine di grandezza di errore):

\begin{center}
\renewcommand{\arraystretch}{1.3}
\begin{tabular}{@{}llccc@{}}
\toprule
\textbf{Fonte} & \textbf{Chip} & $\boldsymbol{p_{01}}$ & $\boldsymbol{p_{10}}$ & \textbf{Rapporto} $p_{10}/p_{01}$\\
\midrule
arXiv:2203.14756 & IBM-Quito, Q0 & $0.0132$ & $0.0468$ & $3.5\times$\\
Documentazione IBM\footnotemark & qubit d'esempio & $0.0295$ & $0.0908$ & $3.1\times$\\
arXiv:2504.12575 & ibmq\_montreal, Q0 & $0.008$ & $0.016$ & $2.0\times$\\
\bottomrule
\end{tabular}
\end{center}
\footnotetext{\url{https://docs.quantum.ibm.com/guides/get-qpu-information}, esempio numerico di backend properties (qubit 126).}

\noindent A questi si aggiunge una conferma qualitativa generale, dalla
stessa documentazione IBM: \emph{``\texttt{prob\_meas0\_prep1} is usually
the dominant factor in the readout error due to relaxation during the
measurement process''} --- cioè $p_{10}>p_{01}$ è la norma su hardware
superconduttore reale, non un'eccezione, per il meccanismo fisico già
discusso nella teoria (rilassamento $T_1$ durante la misura).

\subsection{Scelta dei valori per il test}

\begin{keyeq}[title={Scelta dichiarata}]
\centering
Rapporto $p_{10}:p_{01}=3$ (nel mezzo dell'intervallo $2$--$3.5\times$
osservato sopra), \textbf{stessa media} già in uso $p_\text{readout}=2.3\times10^{-2}$:
$$p_{01}=0.0115, \qquad p_{10}=0.0345$$
\end{keyeq}
Tenere fissa la media ha uno scopo preciso: isolare l'effetto
\emph{dell'asimmetria in sé}, senza mescolarlo con un cambiamento della
scala complessiva del rumore di lettura (che è già stato scansionato
separatamente al Passo 5, nel caso simmetrico). Questi valori sono
\textbf{illustrativi}: non sono lo split reale registrato per
\texttt{ibm\_torino} (non disponibile in forma citabile, Sez.~\ref{sec:ricerca}),
ma un rapporto rappresentativo tratto da calibrazioni IBM reali pubblicate
su chip comparabili --- stesso trattamento già riservato a
$\varepsilon_{2q}$ in \texttt{noise\_model\_dimero.py}.

\section{Verifica}

\begin{okbox}{Quattro controlli superati}
\begin{itemize}[leftmargin=1.2em,itemsep=0.2em]
\item \textbf{Limite di rumore nullo} ($p_{01}=p_{10}=0$): coincide col
ramo originale (readout spento) a precisione macchina ($|\Delta|=0$ su
tutti gli $N$ testati).
\item \textbf{Limite simmetrico} ($p_{01}=p_{10}=p$): riproduce
$N^*=5$ e gli stessi valori $C(N)$ già registrati per il caso simmetrico,
a precisione macchina.
\item \textbf{Cross-check Monte Carlo}: la formula analitica
$\alpha\,C_\text{ideale}+\beta(1+i)$ coincide, entro $5\times10^{-3}$
(compatibile con l'errore statistico atteso a $2\times10^5$ shot), con un
\texttt{ReadoutError} vero e asimmetrico agganciato ad Aer --- due cammini
di calcolo indipendenti per lo stesso risultato.
\item \textbf{Risultato centrale}: $N^*=5$ non si sposta con lo split
illustrativo, né con rapporti fino a $50\times$ (Sez.~\ref{sec:stress}).
\end{itemize}
\end{okbox}

\section{Risultati}

\subsection{Il correlatore rumoroso, simmetrico vs asimmetrico}

\begin{figure}[h]
\centering
\includegraphics[width=0.72\linewidth]{figure/fig_correlatore_N_asimmetrico.pdf}
\caption{\textit{Modulo del correlatore rumoroso $|C_{21}^{xx}(t{=}2,N)|$ in
funzione di $N$, readout simmetrico ($p_{01}=p_{10}=0.023$) contro
asimmetrico ($p_{01}=0.0115$, $p_{10}=0.0345$, stessa media). Il massimo
resta a $N^*=5$ in entrambi i casi. Lo split asimmetrico alza
sistematicamente $|C(N)|$ per $N\gtrsim3$: l'offset $\beta(1+i)$ si somma
qui in modo prevalentemente costruttivo con $C_\text{ideale}(N)$ nella
regione dove il correlatore è già ampio, non è un effetto casuale ma
dipende dalla fase specifica di $C_\text{ideale}(N)$ per questo problema.}}
\label{fig:corr_asimmetrico}
\end{figure}

\begin{center}
\renewcommand{\arraystretch}{1.2}
\begin{tabular}{@{}lcc@{}}
\toprule
\textbf{Configurazione di readout} & $\boldsymbol{N^*}$ & $\boldsymbol{|C(N^*)|}$\\
\midrule
Simmetrico ($p_{01}=p_{10}=0.023$) & $5$ & $0.2500$\\
Asimmetrico illustrativo ($p_{01}=0.0115$, $p_{10}=0.0345$) & $5$ & $0.2775$\\
\bottomrule
\end{tabular}
\end{center}

\subsection{Stress test oltre il range realistico}
\label{sec:stress}

Per capire quanto il risultato $N^*=5$ sia vicino a un limite oppure
robusto, il test è stato ripetuto a media fissa ($2.3\times10^{-2}$) con
rapporti $p_{10}{:}p_{01}$ ben oltre l'intervallo realistico trovato in
letteratura ($2$--$3.5\times$):

\begin{center}
\renewcommand{\arraystretch}{1.2}
\begin{tabular}{@{}cccc@{}}
\toprule
\textbf{Rapporto} $p_{10}{:}p_{01}$ & $\boldsymbol{p_{01}}$ & $\boldsymbol{p_{10}}$ & $\boldsymbol{N^*}$\\
\midrule
$1$ (simmetrico) & $0.0230$ & $0.0230$ & $5$\\
$3$ (illustrativo) & $0.0115$ & $0.0345$ & $5$\\
$5$ & $0.0077$ & $0.0383$ & $5$\\
$10$ & $0.0042$ & $0.0418$ & $5$\\
$20$ & $0.0022$ & $0.0438$ & $5$\\
$50$ & $0.0009$ & $0.0451$ & $5$\\
\bottomrule
\end{tabular}
\end{center}

\begin{figure}[h]
\centering
\includegraphics[width=0.62\linewidth]{figure/fig_Nstar_vs_rapporto.pdf}
\caption{\textit{$N^*$ in funzione del rapporto $p_{10}{:}p_{01}$ (scala
log), a media fissa $2.3\times10^{-2}$. La linea resta piatta a $N^*=5$
per quattro ordini di grandezza di rapporto, ben oltre il range
realistico ($2$--$3.5\times$, Sez.~\ref{sec:ricerca}).}}
\label{fig:Nstar_rapporto}
\end{figure}

\begin{okbox}{$N^*$ empiricamente molto più robusto del previsto dalla teoria}
La teoria (\texttt{teoria\_readout\_asimmetrico.tex}) non garantiva
l'invarianza di $N^*$ per $\beta\neq0$ --- il conto era genuinamente
aperto. Numericamente, per \emph{questo} correlatore del dimero a questo
punto di lavoro, $N^*=5$ risulta invariante non solo nel range realistico
ma fino a un rapporto $50\times$, ben oltre qualunque asimmetria
osservabile su hardware reale.
\end{okbox}

\section{Considerazioni}

\begin{itemize}[leftmargin=1.3em,itemsep=0.35em]
\item \textbf{La domanda era genuinamente aperta, la risposta è
empirica.} A differenza del caso simmetrico (dove l'invarianza di $N^*$ è
un teorema, non serve verificarlo numericamente), qui la teoria lascia
il risultato indeterminato: l'esito --- $N^*$ invariante, e con ampio
margine --- è un fatto numerico specifico di questo correlatore e di
questo punto di lavoro, non una conseguenza necessaria della struttura
del problema. Non è detto che valga per un correlatore diverso, un $t$
diverso, o un altro punto di lavoro $(b/J,D/J)$: è un risultato verificato
qui, non un teorema generale.

\item \textbf{Perché non ha ceduto anche a rapporti estremi.} Guardando
la Fig.~\ref{fig:corr_asimmetrico}, l'offset $\beta(1+i)$ sposta la curva
$C_\text{ideale}(N)$ in una direzione che, per questo correlatore, si
allontana dall'origine più agli $N$ già vicini al picco che agli $N$
lontani da esso --- l'effetto netto è di accentuare il picco esistente
piuttosto che spostarlo. Questo non era garantito dalla teoria (che
ammette anche il caso opposto, offset che appiattisce o sposta il
massimo), ma è coerente con essa: la figura geometrica di
\texttt{teoria\_readout\_asimmetrico.tex} (Fig.~1 di quel documento)
mostrava entrambe le possibilità come geometricamente plausibili.

\item \textbf{Nessun valore reale registrato per \texttt{ibm\_torino}
esiste in forma citabile}, nonostante la ricerca esplicita
(Sez.~\ref{sec:ricerca}). Questo è dichiarato con la stessa onestà con cui
si dichiarano le altre scelte di modellazione della tesi: il risultato
numerico qui presentato usa un rapporto rappresentativo, non lo split
vero del chip citato altrove in tesi per gli altri parametri.

\item \textbf{Criterio di proporzionalità rispettato.} Il relatore ha
chiesto di non investire tempo eccessivo su questo punto opzionale: la
teoria (un documento) più questo test numerico (uno script, una verifica,
questo documento) chiudono la domanda con lo stesso livello di rigore
degli altri passi della pipeline, senza aprire un'indagine sproporzionata
rispetto alla natura facoltativa della richiesta.
\end{itemize}

\section{Conclusioni}

Il test di readout asimmetrico richiesto dal relatore è chiuso: $N^*=5$
(correlatori dinamici, Passo 4) non si sposta introducendo un'asimmetria
$p_{01}\neq p_{10}$, né al punto illustrativo scelto (rapporto $3\times$,
stessa media $2.3\times10^{-2}$) né a rapporti molto più estremi
($50\times$) del range osservato su hardware IBM reale. Questo risultato
non era garantito dalla teoria --- a differenza del caso simmetrico, dove
l'invarianza è dimostrata analiticamente --- ed è quindi un fatto
verificato, non assunto. Nessun valore reale per \texttt{ibm\_torino} è
stato reperito in forma citabile per lo split asimmetrico; i valori usati
sono dichiarati esplicitamente come illustrativi, coerenti con
calibrazioni IBM reali su chip comparabili.

\begin{okbox}{Sintesi in una riga}
Il test opzionale del relatore è chiuso: l'asimmetria del readout non
sposta $N^*$ per questo correlatore, verificato (non solo argomentato)
ben oltre il range realistico.
\end{okbox}

\needspace{6\baselineskip}
\section*{File prodotti}
\addcontentsline{toc}{section}{File prodotti}

\texttt{correlatori\_readout\_asimmetrico.py} (formula generale,
cross-check Monte Carlo, ricerca di $N^*$),
\texttt{validate\_correlatori\_readout\_asimmetrico.py} (i quattro
controlli), questo documento. Vedi anche
\texttt{teoria\_readout\_asimmetrico.tex} per la derivazione completa
usata qui.

\end{document}
